\section{Heterogeneity of the Screening Signal}
\label{sec:mechanisms}

Two pieces of within-evidence discipline the framing. First, the
continuous loss-intensity statistic strictly dominates the
binary rule in discrimination, in line with
Proposition~\ref{prop:exposure_stmt}'s identification of
$\log(1+\text{tenders\_count})$ as the sufficient ranking
statistic (\S\ref{sec:mech_horse_race}). Second, the modal
contrast between preg\~ao and convite reads as variation in the
observability regime within which the screen is recovered
(\S\ref{sec:mech_modal}). Two finer framework predictions---a
cell-heterogeneity grid and a discrete first-time-FL
imprint---are reported in \S\ref{sec:lim_construct} and Online
Appendix~C; the data do not adjudicate them under disciplined
estimation, and the screening interpretation does not require
them.

\subsection{Loss Intensity Dominates the Threshold Cutoff}
\label{sec:mech_horse_race}

The binary rule is the information-coarsening of the
participation primitive identified by the framework, not a
primitive in itself. On a harmonized same-sample basis, the
continuous $\log(1+\text{tenders\_count})$ score discriminates
the $\valCobidders$ adjudicated CADE-cobidder firms inside the
always-loser stratum at AUC $\valAUClogtc$
($95\%$ CI $\valAUClogtcCI$) versus AUC $\valAUCFLBinSameSample$
($\valAUCFLBinSameSampleCI$) for the binary FL14 indicator
(\citealt{delong1988comparing} test, $Z = \valDeLongZ$,
$p = \valDeLongP$): the $\sim\!2.8$-point AUC gain on the same
firm panel is the dose-response signature the framework predicts
when the binary coarsens a continuous primitive. Standalone price
coefficients align (Table~\ref{tab:horse_race_v14}, columns
1--2): FL14 binary $+0.0653$ ($p<0.001$), continuous
$\log(1+\max\,\text{tc})$ $+0.0188$ ($p<0.001$). When both are
included additively (column 3), the FL14 coefficient flips sign
($-0.0746$, $p<0.05$) while the continuous coefficient nearly
doubles ($+0.0349$, $p<0.001$): once the continuous primitive is
in the specification, the binary indicator captures only tail
residuals, and reading the binary in isolation overstates the
cover-bidder type information by absorbing the variation the
continuous primitive carries. The dose-response is consistent
with the screening interpretation and inconsistent with a
pure-confound reading of the broad-sample association.

% src: scripts/34_horse_race_fl_continuous.R + 36_gate_d1_harmonized.R
\begin{table}[htbp]
\centering
\caption{Horse Race: FL14 Binary vs Continuous Loss Intensity (Same Sample)}
\label{tab:horse_race_v14}
\begin{threeparttable}
\small
\begin{tabular}{lcccc}
\toprule
Specification & Coef on FL14 & Coef on $\log(1{+}\max\,\text{tc})$ & Convite & $N$ \\
\midrule
(1) FL14 only                 & $\valHorseFLOne$ & --                & $\valHorseConvOne$ & $\valHorseN$ \\
                              & $(\valHorseFLOneSE)$  &              & $(\valHorseConvOneSE)$ & \\
(2) Continuous only           & --                & $\valHorseContTwo$ & $\valHorseConvTwo$ & $\valHorseN$ \\
                              &                   & $(\valHorseContTwoSE)$ & $(\valHorseConvTwoSE)$ & \\
(3) Both, additive            & $\valHorseFLThree$ & $\valHorseContThree$ & $\valHorseConvThree$ & $\valHorseN$ \\
                              & $(\valHorseFLThreeSE)$ & $(\valHorseContThreeSE)$ & $(\valHorseConvThreeSE)$ & \\
\midrule
\multicolumn{5}{l}{\emph{Firm-level AUC against $\valCobidders$ CADE-cobidder labels (always-loser pool, $N=\valAlwaysLosers$)}} \\
FL14 binary                       & AUC $= \valHorseAUCBin$ $\valHorseAUCBinCI$ & & & \\
$\log(1{+}\text{tenders\_count})$ & AUC $= \valHorseAUCCont$ $\valHorseAUCContCI$ & & & \\
DeLong test (continuous vs FL14)  & $Z = \valDeLongZ$, $p = \valDeLongP$ & & & \\
\bottomrule
\end{tabular}
\begin{tablenotes}
\footnotesize
\item \textit{Notes:} Outcome is $\log p_{\text{negotiated}}$ within the
v13 panel of $\valHorseN$ items. Item-level treatment is
$\mathbf{1}[\text{any FL14 participant}]$ in column~(1) and the log of the
maximum loss-intensity score among always-loser participants in
column~(2). Column~(3) includes both. The FL14 coefficient flips sign
when continuous intensity is added: the binary discretization at $14$
tenders is the information-coarsening of the underlying continuous
signal, not an economically primitive threshold. SEs clustered at item
level. *** $p<0.01$, ** $p<0.05$, * $p<0.1$.
\end{tablenotes}
\end{threeparttable}
\end{table}


We use the binary FL14 rule as the discrete cut-off and the
continuous statistic as the discriminating instrument: the
headline range $\valHeadlineRange$ in \S\ref{sec:results_ols} is
computed on the binary, and the AUC headline $\valAUCFLfirm$
($95\%$ CI $\valAUCFLfirmCI$) reported in
\S\ref{sec:cade_prospective} is binary firm-level discrimination.

\subsection{Modal Contrast as Observability-Regime Variation}
\label{sec:mech_modal}

The screening signal is sharper in preg\~ao than in convite
along two independent dimensions. On the price-imprint side, the
within-modal estimates of \S\ref{sec:results_ols} place the
binary specification at $\valFalPregBin$ ($\valFalPregBinPSig$)
inside preg\~ao versus $\valFalConvBin$ ($\valFalConvBinPSig$)
inside convite, a ratio of $\valFalRatio\!\!\times$ on the larger
preg\~ao sample. On the discrimination side, the modal-by-modal
AUC against the $\valCobidders$ adjudicated CADE-cobidder firms
is $\valAUCpregPrim$ in preg\~ao primary auctions versus
$\valAUCconvPrim$ in convite primary auctions, with a bootstrap
difference of $\valAUCmodalDiff$ ($p\approx 0$). The two
dimensions point in the same direction; the convite-modal AUC
carries a sample-size caveat ($\valConvNcobidders$ adjudicated
cobidders inside the convite stratum against several dozen on
the preg\~ao side) and we treat it as a directional indicator,
with the much larger $\valFalNConv$-tender-item price-imprint
evidence carrying the statistical weight.

The convite minimum-bidder rule
(Lei~$8.666/93$, Art.~$22$ \S\,$3$) is the obvious institutional
suspect: a quorum-filler reading would predict a sharper
screening signal where the rule binds, that is, in convite. The
data reverse this prediction. We do not read the reversal as a
positive test of any alternative institutional theory, but the
contrast is consistent with a structural feature of the two
modalities that the cover-bidding literature has long
recognized. Cover bidding in a sealed-bid environment is a
single \emph{passive} act: the cover bidder submits an envelope
and the participation footprint is one observation per item.
Cover bidding in a real-time electronic auction is a sequence
of \emph{active} acts: the cover bidder must respond to other
bids, maintain relevance without winning, and replicate
plausible timing patterns over the auction's lifespan. Each
revision is an additional observation against which the
deployment pattern is read. Pregão therefore does not just
expose a richer interaction footprint than convite---it forces
the cover bidder to leave more of one. This is the
information-theoretic counterpart of the dynamic-vs-sealed-bid
distinction in
\citet{marshall2012economics,asker2010leniency,athey2011comparing}:
a sealed envelope hides whatever the auction does not call out
of it; a real-time auction calls more out. Two compatible
mechanisms operate alongside this---bidder-pool composition
that survives convite's quorum requirement may itself be more
homogeneous than the pregão pool, and the observability regime
shapes what the analyst can read off the data envelope---and
the design separates none of the three from each other. What
the design \emph{does} establish is the asymmetry's direction
and its alignment with how the cover-bidder type behaves under
each modality. We therefore report the modal asymmetry as
\emph{scope information} for the screening object---the
construct discriminates better in pregão environments---and
not as a positive test of the minimum-bidder-rule mechanism,
which would require institutional variation the BEC setting
does not deliver.

The implication for portability is concrete. The post-sample
reform (Lei~$14.133/2021$) consolidates pregão-style auctions as
the institutional default, which is the regime in which the
screening signal is sharpest. The construct travels best to
jurisdictions whose observability regime resembles
pregão---real-time electronic auctions with public
bidder-by-bidder interaction---and the regulatory direction is
favourable rather than adversarial.
